\documentclass[10pt,a4paper]{amsart}
\usepackage[margin=19mm]{geometry}
\usepackage{amsmath,amssymb,mathtools}
\usepackage[scr=rsfs,cal=euler]{mathalfa}
\usepackage{xfrac}
\usepackage[unicode,hidelinks,bookmarks=false]{hyperref}
\newcommand{\ie}{i.e.\ }
\newcommand{\cf}{cf.\ }
\newcommand{\resp}{respectively\ }
\newcommand{\Subsection}{\subsection}
\providecommand{\nobreakdash}{\nobreak\hbox{-}}
\setlength{\emergencystretch}{2em}
\multlinegap0pt
\usepackage{xcolor}
\usepackage{amssymb}
\usepackage{tikz}
\usepackage{enumerate}
\usepackage{float}
\newtheorem{prop}{Proposition}
\newtheorem{lem}{Lemma}
\title{Subrack lattices of finite solvable and metacyclic groups}
\author{Sel\c{c}uk Kayacan}
\date{Source: arXiv:2503.06714v2 (CC BY 4.0)}
\begin{document}
\maketitle

\noindent\emph{Source and licence.} Subrack lattices of finite solvable and metacyclic groups. Version arXiv:2503.06714v2 (CC BY 4.0), distributed by arXiv under the Creative Commons Attribution 4.0 International licence, http://creativecommons.org/licenses/by/4.0/, as recorded in the arXiv licence metadata for this exact version and shown on the abstract page https://arxiv.org/abs/2503.06714v2. Full licence notice: \texttt{licenses/arxiv-2503.06714-CC-BY-4.0.txt}.\par\smallskip

\noindent\textit{Editorial scope.} Counted unit: the article's prop prop:ex (source0.tex lines 164--171). The article's complete original proof of it is delivered with it (source0.tex lines 664--860, one \texttt{proof} environment of 197 lines, which the article places away from the statement). No mathematics is written, repaired or re-derived: the statement and the proof are the article's own lines, in the article's own notation, and no retained line is substituted or re-flowed.  Retained uncounted as the source-local context the proof needs: lines 566--571; lines 600--602; lines 648--651.  Nothing was omitted from any retained span, so the package carries no omission record.  No retained line contains a citation, so the wrapper carries no bibliography and no bibliography entry.  No retained line contains a figure environment, an image-inclusion command or drawing code, so no image is delivered and the package declares no asset dependency.  Editorial formatting only, in addition: an AMS \texttt{amsart} preamble that restates the article's own declarations and macros used by the retained lines, namely the environments \texttt{prop}, \texttt{lem} on separate counters, together with the standard packages those lines need.  The retained environments print in this document as Proposition 1, Lemma 1 in order of appearance, whereas the article prints the same items with the numbering of its own full document.  The complete native source of the article as distributed in the arXiv e-print for this exact version is bundled unchanged as \texttt{sources/174462/source0.tex}.
\begin{prop}\label{prop:pairs}
  Let $X$ and $Y$ be two finite racks, and let $f\colon X\to Y$ be a bijection satisfying
  $$ f(\langle x,x'\rangle_{\mathsf{rk}}) = \langle f(x),f(x')\rangle_{\mathsf{rk}} $$
  for all $x,x'\in X$. Then $f$ induces a lattice isomorphism
  $$ f\colon\mathcal{R}(X)\to \mathcal{R}(Y). $$
\end{prop}

\begin{lem}\label{lem:g0h0}
  The groups $G_0$ and $H_0$ are not isomorphic.
\end{lem}

\begin{lem}\label{lem:pair}
  Let $G$ be a group and let $x,y\in G$. Set $L:=\langle x,y\rangle$. Then
  $$ \langle x,y\rangle_{\mathsf{rk}}=K_L(x)\cup K_L(y). $$
\end{lem}

\begin{prop}\label{prop:ex}
  Let
  \begin{align*}
    G_0 &= \langle a,b\mid a^{91} = 1,\ b^3 = 1,\ bab^{-1} = a^9\rangle, \\
    H_0 &= \langle \alpha,\beta\mid \alpha^{91} = 1,\ \beta^3 = 1,\ \beta\alpha\beta^{-1} = \alpha^{16}\rangle.
  \end{align*}
  Then $\mathcal{R}(G_0)\cong \mathcal{R}(H_0)$ but $G_0\not\cong_{\mathsf{rk}}H_0$.
\end{prop}

\begin{proof}[Proof of Proposition~\ref{prop:ex}]
  First we show that $G_0$ and $H_0$ are not isomorphic as racks. Since $Z(G_0)=Z(H_0)=1$, any isomorphism between the group racks $G_0$ and $H_0$ would induce an isomorphism of groups $G_0\to H_0$. By Lemma~\ref{lem:g0h0}, the groups $G_0$ and $H_0$ are not isomorphic. Therefore $G_0\not\cong_{\mathsf{rk}} H_0$.

  Next we show that the map $\psi_0\colon G_0\to H_0$ induces a lattice isomorphism. We verify the equality
  $$ \psi_0(\langle x,y\rangle_{\mathsf{rk}}) = \langle \psi_0(x),\psi_0(y)\rangle_{\mathsf{rk}} $$
  for every pair $x,y\in G_0$. Since every element of $G_0$ lies in exactly one of the three cosets $N=\langle a\rangle$, $Nb$, and $Nb^2$, every pair falls into exactly one of the six unordered coset-types
  $$ (N,N),\ (N,Nb),\ (N,Nb^2),\ (Nb,Nb),\ (Nb,Nb^2),\ (Nb^2,Nb^2). $$
  In each nonabelian case the argument follows the same pattern: we set $L:=\langle x,y\rangle$, determine $L\cap N=\langle a^d\rangle$ for a suitable divisor $d$ of $91$, observe that $(a^d,x)$ or $(a^d,y)$ is a Hölder pair for $L$, describe the relevant conjugacy classes as $\langle a^d\rangle$-orbits, and then apply Lemma~\ref{lem:pair}. We then repeat the same analysis for $L':=\langle \psi_0(x),\psi_0(y)\rangle$, using that $\psi_0$ preserves the order classes $A_1,A_7,A_{13},A_{91}$ inside $N$ and maps each conjugacy class in $N$ onto the corresponding conjugacy class in $\langle \alpha\rangle$.

  \medskip
  \noindent
  \emph{Case~1.\; $\psi_0(\langle a^{i_1},a^{i_2}\rangle_{\mathsf{rk}}) = \langle \psi_0(a^{i_1}),\psi_0(a^{i_2})\rangle_{\mathsf{rk}}$}

  \medskip
  Let $x=a^{i_1}$ and $y=a^{i_2}$. Since $N$ and $\psi_0(N)=\langle\alpha\rangle$ are abelian, both $\langle x,y\rangle_{\mathsf{rk}}$ and $\langle \psi_0(x),\psi_0(y)\rangle_{\mathsf{rk}}$ are just the subracks generated by two commuting elements. Hence the desired equality follows.
  
  
  \medskip
  \noindent
  \emph{Case~2.\; $\psi_0(\langle a^{i_1},a^{i_2}b\rangle_{\mathsf{rk}})=\langle \psi_0(a^{i_1}),\psi_0(a^{i_2}b)\rangle_{\mathsf{rk}}$}
  
  \medskip
  If $i_1=0$, then $a^{i_1}=1$ is central, so $\langle a^{i_1},a^{i_2}b\rangle_{\mathsf{rk}}=\{1,a^{i_2}b\}$. Since $\psi_0(1)=1$ and $\psi_0(a^{i_2}b)=\alpha^{i_2}\beta$, we also have $\langle \psi_0(a^{i_1}),\psi_0(a^{i_2}b)\rangle_{\mathsf{rk}}=\{1,\alpha^{i_2}\beta\}$, and the desired equality follows.
  
  Assume now that $i_1\neq 0$. Set
  $$ x:=a^{i_1},\qquad y:=a^{i_2}b,\qquad L:=\langle x,y\rangle. $$
  Let $d:=\gcd(91,i_1)$ and $m:=91/d$. Then $|x|=m$, and
  $$ yxy^{-1} = (a^{i_2}b)a^{i_1}(a^{i_2}b)^{-1} = ba^{i_1}b^{-1} = a^{9i_1} = x^9. $$
  Also, $(a^{i_2}b)^3=1$. Hence $(x,y)$ is a Hölder pair for $L$ with $L=\langle x,y\mid m,3,0,9\rangle$. Since $(x,y)$ is a Hölder pair for $L$, the conjugacy class of $x$ in $L$ is the
  $\langle y\rangle$-orbit of $x$, namely
  $$ K_L(x)=\{x,x^9,x^{81}\}=\{a^{i_1},a^{9i_1},a^{81i_1}\}. $$
  Similarly, the conjugacy class of $y$ in $L$ is the $\langle x\rangle$-orbit of $y$. Since $x^n y x^{-n}=a^{i_2-8n i_1}b$ and $-8$ is a unit modulo $91$, we obtain
  $$ K_L(y)=\{a^j b : j\in i_2+\langle i_1\rangle\}, $$
  where $\langle i_1\rangle$ denotes the additive subgroup of $\mathbb{Z}/91\mathbb{Z}$ generated by $i_1$. Hence, by Lemma~\ref{lem:pair},
  $$ \langle a^{i_1},a^{i_2}b\rangle_{\mathsf{rk}} = K_L(x)\cup K_L(y) = K_{G_0}(a^{i_1})\cup\{a^j b : j\in i_2+\langle i_1\rangle\}. $$
  
  Now set
  $$ x':=\psi_0(a^{i_1}),\qquad y':=\psi_0(a^{i_2}b)=\alpha^{i_2}\beta,\qquad L':=\langle x',y'\rangle. $$
  Since $\psi_0$ preserves the order of elements of $\langle a\rangle$, we have $|x'|=|x|=m$. Moreover,
  $$ y'x'(y')^{-1}=\beta x'\beta^{-1}=(x')^{16}, $$
  and $(y')^3=1$. Thus $(x',y')$ is a Hölder pair for $L'$ with $L'=\langle x',y'\mid m,3,0,16\rangle$. Since $\psi_0$ preserves the order of $a^{i_1}$ and hence the additive subgroup $\langle i_1\rangle\leq \mathbb{Z}/91\mathbb{Z}$, we have
  $$ K_{L'}(x')=K_{H_0}(x')\quad \text{and}\quad K_{L'}(y')=\{\alpha^j\beta : j\in i_2+\langle i_1\rangle\}. $$
  Hence, by Lemma~\ref{lem:pair},
  $$ \langle x',y'\rangle_{\mathsf{rk}} = K_{L'}(x')\cup K_{L'}(y') = K_{H_0}(\psi_0(a^{i_1}))\cup \{\alpha^j\beta : j\in i_2+\langle i_1\rangle\}. $$
  
  
  
  \medskip
  \noindent
  \emph{Case~3.\; $\psi_0(\langle a^{i_1},a^{i_2}b^2\rangle_{\mathsf{rk}})=\langle \psi_0(a^{i_1}),\psi_0(a^{i_2}b^2)\rangle_{\mathsf{rk}}$}
  
  \medskip
  If $i_1=0$, then $a^{i_1}=1$ is central, so $\langle a^{i_1},a^{i_2}b^2\rangle_{\mathsf{rk}}=\{1,a^{i_2}b^2\}$. Since $\psi_0(1)=1$ and $\psi_0(a^{i_2}b^2)=\alpha^{29i_2}\beta^2$, we also have $\langle \psi_0(a^{i_1}),\psi_0(a^{i_2}b^2)\rangle_{\mathsf{rk}}=\{1,\alpha^{29i_2}\beta^2\}$, and the desired equality follows.
  
  Assume now that $i_1\neq 0$. Set
  $$ x:=a^{i_1},\qquad y:=a^{i_2}b^2,\qquad L:=\langle x,y\rangle. $$
  Let $d:=\gcd(91,i_1)$ and $m:=91/d$. Then $|x|=m$, and
  $$ yxy^{-1} = (a^{i_2}b^2)a^{i_1}(a^{i_2}b^2)^{-1} = b^2a^{i_1}b^{-2} = a^{81i_1} = x^{81}. $$
  Also, $(a^{i_2}b^2)^3=1$. Hence $(x,y)$ is a Hölder pair for $L$ with $L=\langle x,y\mid m,3,0,81\rangle$ and the conjugacy class of $x$ in $L$ is the $\langle y\rangle$-orbit of $x$, namely
  $$ K_L(x)=\{x,x^{81},x^{81^2}\}=\{a^{i_1},a^{81i_1},a^{9i_1}\}=K_{G_0}(a^{i_1}). $$
  Similarly, the conjugacy class of $y$ in $L$ is the $\langle x\rangle$-orbit of $y$. Since
  $$ x^n y x^{-n} = a^{ni_1}(a^{i_2}b^2)a^{-ni_1} = a^{\,i_2-80ni_1}b^2, $$
  and $-80$ is a unit modulo $91$, we obtain
  $$ K_L(y)=\{a^j b^2 : j\in i_2+\langle i_1\rangle\}, $$
  where $\langle i_1\rangle$ denotes the additive subgroup of $\mathbb{Z}/91\mathbb{Z}$ generated by $i_1$. Hence, by Lemma~\ref{lem:pair},
  $$ \langle a^{i_1},a^{i_2}b^2\rangle_{\mathsf{rk}} = K_L(x)\cup K_L(y) = K_{G_0}(a^{i_1})\cup \{a^j b^2 : j\in i_2+\langle i_1\rangle\}. $$
  
  Now set
  $$ x':=\psi_0(a^{i_1}),\qquad y':=\psi_0(a^{i_2}b^2)=\alpha^{29i_2}\beta^2,\qquad L':=\langle x',y'\rangle. $$
  Since $\psi_0$ preserves the order of elements of $\langle a\rangle$, we have $|x'|=|x|=m$. Also,
  $$ y'x'(y')^{-1} = (\alpha^{29i_2}\beta^2)x'(\alpha^{29i_2}\beta^2)^{-1} = \beta^2x'\beta^{-2} = (x')^{74}, $$
  because if $x'=\alpha^u$, then $\beta^2\alpha^u\beta^{-2}=\alpha^{16^2u}=\alpha^{74u}$. Moreover, $(y')^3=1$. Hence $(x',y')$ is a Hölder pair for $L'$ with $L' = \langle x',y'\mid m,3,0,74\rangle$. Therefore the conjugacy class of $x'$ in $L'$ is the $\langle y'\rangle$-orbit of $x'$,
  that is,
  $$ K_{L'}(x')=K_{H_0}(x'). $$
  Similarly, the conjugacy class of $y'$ in $L'$ is the $\langle x'\rangle$-orbit of $y'$. If
  $x'=\alpha^u$, then
  $$ (x')^n y' (x')^{-n} = \alpha^{nu}(\alpha^{29i_2}\beta^2)\alpha^{-nu} = \alpha^{\,29i_2-73nu}\beta^2. $$
  Since $\psi_0$ preserves the gcd-layer of $i_1$, the additive subgroups generated by $u$, $i_1$, and $29i_1$ in $\mathbb{Z}/91\mathbb{Z}$ all coincide. As $-73$ is a unit modulo
  $91$, it follows that
  $$ K_{L'}(y')=\{\alpha^k\beta^2 : k\in 29i_2+\langle 29i_1\rangle\}. $$
  On the other hand, $\psi_0(a^j b^2)=\alpha^{29j}\beta^2$, so
  \begin{align*}
    \psi_0(\{a^j b^2:j\in i_2+\langle i_1\rangle\})
    &= \{\alpha^{29j}\beta^2:j\in i_2+\langle i_1\rangle\} \\
    &= \{\alpha^k\beta^2:k\in 29i_2+\langle 29i_1\rangle\}.
  \end{align*}
  Thus
  $$ \psi_0(K_L(y))=K_{L'}(y') \qquad\text{and}\qquad \psi_0(K_L(x))=K_{L'}(x'). $$
  Hence, by Lemma~\ref{lem:pair},
  $$ \psi_0(\langle a^{i_1},a^{i_2}b^2\rangle_{\mathsf{rk}}) = K_{L'}(x')\cup K_{L'}(y') = \langle \psi_0(a^{i_1}),\psi_0(a^{i_2}b^2)\rangle_{\mathsf{rk}}. $$
  
  
  
  \medskip
  \noindent
  \emph{Case~4.\; $\psi_0(\langle a^{i_1}b,a^{i_2}b\rangle_{\mathsf{rk}}) = \langle \psi_0(a^{i_1}b),\psi_0(a^{i_2}b)\rangle_{\mathsf{rk}}$}
  
  \medskip
  Set
  $$ x:=a^{i_1}b,\qquad y:=a^{i_2}b,\qquad L:=\langle x,y\rangle. $$
  Let $d:=\gcd(91,i_1,i_2)$. Then $L\cap \langle a\rangle = \langle a^d\rangle$ and $(a^d,x)$ is a Hölder pair for $L$ with $L/\langle a^d\rangle$ is cyclic of order $3$. Therefore the conjugacy class of $x$ in $L$ is the $\langle a^d\rangle$-orbit of $x$. Since
  $$ a^{nd}xa^{-nd}=a^{nd}(a^{i_1}b)a^{-nd}=a^{\,i_1-8nd}b, $$
  and $-8$ is a unit modulo $91$, we obtain
  $$ K_L(x)=\{a^{i_1+t}b : t\in \langle d\rangle\}, $$
  where $\langle d\rangle$ denotes the additive subgroup of $\mathbb{Z}/91\mathbb{Z}$ generated by $d$. Since $x$ and $y$ are conjugate in $L$, by Lemma~\ref{lem:pair} we have
  $$ \langle a^{i_1}b,a^{i_2}b\rangle_{\mathsf{rk}} = K_L(x)\cup K_L(y) = \{a^{i_1+t}b : t\in \langle d\rangle\}. $$
  
  By construction,
  $$ \psi_0(a^jb)=\alpha^j\beta \qquad (0\leq j<91), $$
  so
  $$ \psi_0(\{a^{i_1+t}b : t\in \langle d\rangle\}) = \{\alpha^{i_1+t}\beta : t\in \langle d\rangle\}. $$
  Now set
  $$ x':=\alpha^{i_1}\beta,\qquad y':=\alpha^{i_2}\beta,\qquad L':=\langle x',y'\rangle. $$
  As above $L'\cap \langle \alpha\rangle=\langle \alpha^d\rangle$ and $(\alpha^d,x')$ is a Hölder pair for $L'$ with $L'/\langle\alpha^d \rangle$ is cyclic of order $3$. Therefore the conjugacy class of $x'$ in $L'$ is the $\langle \alpha^d\rangle$-orbit of $x'$. Since
  $$ \alpha^{nd}x'\alpha^{-nd} = \alpha^{nd}(\alpha^{i_1}\beta)\alpha^{-nd} = \alpha^{\,i_1-15nd}\beta, $$
  and $-15$ is a unit modulo $91$, we obtain
  $$ K_{L'}(x')=\{\alpha^{i_1+t}\beta : t\in \langle d\rangle\}. $$
  Since $x'$ and $y'$ are conjugate in $L'$, by Lemma~\ref{lem:pair} we have 
  $$ \langle x',y'\rangle_{\mathsf{rk}} = K_{L'}(x')\cup K_{L'}(y') = \{\alpha^{i_1+t}\beta : t\in \langle d\rangle\}. $$
  Therefore
  $$ \psi_0(\langle a^{i_1}b,a^{i_2}b\rangle_{\mathsf{rk}}) = \langle \alpha^{i_1}\beta,\alpha^{i_2}\beta\rangle_{\mathsf{rk}} = \langle \psi_0(a^{i_1}b),\psi_0(a^{i_2}b)\rangle_{\mathsf{rk}}. $$
  
  
  \medskip
  \noindent
  \emph{Case~5.\; $\psi_0(\langle a^{i_1}b,a^{i_2}b^2\rangle_{\mathsf{rk}}) = \langle \psi_0(a^{i_1}b),\psi_0(a^{i_2}b^2)\rangle_{\mathsf{rk}}$}
  
  \medskip
  Set
  $$ x:=a^{i_1}b,\qquad y:=a^{i_2}b^2,\qquad L:=\langle x,y\rangle. $$
  Let $d:=\gcd(91,i_1,i_2)$. Then $L\cap \langle a\rangle = \langle a^d\rangle$ and $(a^d,x)$ is a Hölder pair for $L$ with $L/\langle a^d\rangle$ is cyclic of order $3$. Therefore the conjugacy class of $x$ in $L$ is the $\langle a^d\rangle$-orbit of $x$. Since
  $$ a^{nd}xa^{-nd}=a^{nd}(a^{i_1}b)a^{-nd}=a^{\,i_1-8nd}b, $$
  and $-8$ is a unit modulo $91$, we obtain
  $$ K_L(x)=\{a^{i_1+t}b : t\in \langle d\rangle\}, $$
  where $\langle d\rangle$ denotes the additive subgroup of $\mathbb{Z}/91\mathbb{Z}$
  generated by $d$. Similarly, the conjugacy class of $y$ in $L$ is the $\langle a^d\rangle$-orbit of $y$. Since
  $$ a^{nd}ya^{-nd}=a^{nd}(a^{i_2}b^2)a^{-nd}=a^{\,i_2-80nd}b^2, $$
  and $-80$ is a unit modulo $91$, we obtain
  $$ K_L(y)=\{a^{i_2+t}b^2 : t\in \langle d\rangle\}. $$
  Hence, by Lemma~\ref{lem:pair},
  $$ \langle a^{i_1}b,a^{i_2}b^2\rangle_{\mathsf{rk}} = K_L(x)\cup K_L(y) = \{a^{i_1+t}b : t\in \langle d\rangle\} \,\cup\, \{a^{i_2+t}b^2 : t\in \langle d\rangle\}. $$
  
  Now set
  $$ x':=\psi_0(x)=\alpha^{i_1}\beta,\qquad y':=\psi_0(y)=\alpha^{29i_2}\beta^2,\qquad L':=\langle x',y'\rangle. $$
  As above $L'\cap \langle \alpha\rangle=\langle \alpha^d\rangle$ and $(\alpha^d,x')$ is a Hölder pair for $L'$ with $L'/\langle\alpha^d \rangle$ is cyclic of order $3$. Since the conjugacy class of $x'$ in $L'$ is the $\langle \alpha^d\rangle$-orbit of $x'$, we obtain
  $$
  K_{L'}(x')=\{\alpha^{i_1+t}\beta : t\in \langle d\rangle\}.
  $$
  Similarly, the conjugacy class of $y'$ in $L'$ is the $\langle \alpha^d\rangle$-orbit of $y'$ and we have
  $$ K_{L'}(y')=\{\alpha^{29i_2+t}\beta^2 : t\in \langle d\rangle\}. $$
  
  By construction,
  $$ \psi_0(a^{i_1+t}b)=\alpha^{i_1+t}\beta \qquad\text{and}\qquad \psi_0(a^{i_2+t}b^2)=\alpha^{29(i_2+t)}\beta^2. $$
  Since multiplication by $29$ is a unit modulo $91$, the sets
  $$ \{\alpha^{29(i_2+t)}\beta^2:t\in \langle d\rangle\} \qquad\text{and}\qquad \{\alpha^{29i_2+t}\beta^2:t\in \langle d\rangle\} $$
  coincide. Hence
  $$ \psi_0(K_L(x))=K_{L'}(x') \qquad\text{and}\qquad \psi_0(K_L(y))=K_{L'}(y'). $$
  Therefore, again by Lemma~\ref{lem:pair},
  $$ \psi_0(\langle a^{i_1}b,a^{i_2}b^2\rangle_{\mathsf{rk}}) = \langle \alpha^{i_1}\beta,\alpha^{29i_2}\beta^2\rangle_{\mathsf{rk}} = \langle \psi_0(a^{i_1}b),\psi_0(a^{i_2}b^2)\rangle_{\mathsf{rk}}. $$
  
  
  
  \medskip
  \noindent
  \emph{Case~6.\; $\psi_0(\langle a^{i_1}b^2,a^{i_2}b^2\rangle_{\mathsf{rk}})
    =\langle \psi_0(a^{i_1}b^2),\psi_0(a^{i_2}b^2)\rangle_{\mathsf{rk}}$}
  
  \medskip
  Set
  $$ x:=a^{i_1}b^2,\qquad y:=a^{i_2}b^2,\qquad L:=\langle x,y\rangle. $$
  Let $d:=\gcd(91,i_1,i_2)$. Similar to previous cases $(a^d,x)$ is a Hölder pair for $L$ and the conjugacy class of $x$ in $L$ is the $\langle a^d\rangle$-orbit of $x$. Since
  $$ a^{nd}xa^{-nd}=a^{nd}(a^{i_1}b^2)a^{-nd}=a^{\,i_1-80nd}b^2, $$
  and $-80$ is a unit modulo $91$, we have
  $$ K_L(x)=\{a^{i_1+t}b^2:t\in \langle d\rangle\}. $$
  Since $x$ and $y$ are conjugate in $L$, by Lemma~\ref{lem:pair} we have
  $$ \langle a^{i_1}b^2,a^{i_2}b^2\rangle_{\mathsf{rk}} = K_L(x)\cup K_L(y) = \{a^{i_1+t}b^2 : t\in \langle d\rangle\}. $$
  
  As was observed before
  $$ \psi_0(\{a^{i_1+t}b^2:t\in \langle d\rangle\}) = \{\alpha^{29(i_1+t)}\beta^2:t\in \langle d\rangle\} = \{\alpha^{29i_1+t}\beta^2:t\in \langle d\rangle\}. $$
  Now set
  $$ x':=\alpha^{29i_1}\beta^2,\qquad y':=\alpha^{29i_2}\beta^2,\qquad L':=\langle x',y'\rangle. $$
  Similar to Case~4, $(\alpha^d,x')$ is a Hölder pair for $L'$ and the conjugacy class of $x'$ in $L'$ is the $\langle \alpha^d\rangle$-orbit of $x'$. Since
  $$ \alpha^{nd}x'\alpha^{-nd} = \alpha^{nd}(\alpha^{29i_1}\beta^2)\alpha^{-nd} = \alpha^{\,29i_1-73nd}\beta^2, $$
  and $-73$ is a unit modulo $91$, we obtain
  $$ K_{L'}(x')=\{\alpha^{29i_1+t}\beta^2:t\in \langle d\rangle\}. $$
  Since $x'$ and $y'$ are conjugate in $L'$, by Lemma~\ref{lem:pair} we have 
  $$ \langle x',y'\rangle_{\mathsf{rk}} = K_{L'}(x')\cup K_{L'}(y') = \{\alpha^{29i_1+t}\beta^2:t\in \langle d\rangle\}. $$
  Therefore
  $$ \psi_0(\langle a^{i_1}b^2,a^{i_2}b^2\rangle_{\mathsf{rk}}) = \langle \alpha^{29i_1}\beta^2,\alpha^{29i_2}\beta^2\rangle_{\mathsf{rk}} = \langle \psi_0(a^{i_1}b^2),\psi_0(a^{i_2}b^2)\rangle_{\mathsf{rk}}. $$
  
  
  
  \medskip
  The six cases above exhaust all possibilities for pairs of elements of $G_0$. Hence $\psi_0$ satisfies the hypothesis of Proposition~\ref{prop:pairs}. Therefore $\psi_0$ induces a lattice isomorphism from $\mathcal{R}(G_0)$ to $\mathcal{R}(H_0)$. Combining this with the first part of the proof, we obtain Proposition~\ref{prop:ex}.
  
  
\end{proof}

\end{document}
